% Source: https://artofproblemsolving.com/community/c3h3166099_2019_lmt_spring_guts_round_p25p36__lexington_mathematical_tournament
% Stable identity: aops-topic-3166099-q25
Place $A=(-1,0)$, $B=(1,0)$, and the unit-circle center at the origin. If the smaller circle has radius $r$ and center $(h,r)$, internal tangency gives $\sqrt{h^2+r^2}=1-r$, implying $r\le\dfrac{1}{2}$, with equality only at $h=0$. Thus the incircle has center $(0,\dfrac{1}{2})$ and radius $\dfrac{1}{2}$. The nonhorizontal tangent from $A$ has equation $y=m(x+1)$. Its distance from the center must be $\dfrac{1}{2}$, so $(m-\dfrac{1}{2})^2=\dfrac{1}{4}(m^2+1)$, yielding $m=\dfrac{4}{3}$ apart from the base tangent. By symmetry the other side has slope $-\dfrac{4}{3}$, meeting it at $C=(0,\dfrac{4}{3})$. The area is $\dfrac{1}{2}\cdot2\cdot\dfrac{4}{3}=\boxed{\dfrac{4}{3}}$.
\begin{asy}
usepackage("amsmath");
size(7cm);
pair A=(-1,0),Bp=(1,0),C=(0,4/3);draw(circle((0,0),1),gray);draw(circle((0,.5),.5),blue);draw(A--Bp--C--cycle);label("$A$",A,W);label("$B$",Bp,E);label("$C$",C,N);
\end{asy}
